%%%% Introduction
Software is provided to users through HTTP interfaces that sit on opaque stacks of virtualised compute, managed platforms, third-party APIs, and geographically distributed infrastructure. This architecture has improved reliability and development velocity but it has made environmental impact harder to observe and harder to attribute \cite{parssinen2018onlineadvertising,han2025fairco2}. Users can see latency and price, yet they usually cannot see the energy use or emissions associated with an individual action. Providers face the complementary problem: even when they want to disclose environmental impact, request handling often spans caches, databases, queues, and external services, while direct power telemetry is unavailable or too coarse in cloud environments.

A recurring principle in sustainable computing is that information should be exposed at the interface where decisions are actually made. The Software Carbon Intensity (SCI) specification reflects this by expressing emissions relative to a functional unit \(R\), rather than only as annual or service-level aggregates \cite{gsf-sci}. For web systems, the natural functional unit is the individual HTTP transaction. If request-level impact were available at the protocol boundary, clients could directly incorporate it into software behaviour, for example in retry strategies, batching, caching, provider selection, or prompt design.

HTTP is a natural candidate for this disclosure as it already serves as the interoperability layer across services, organizations, and infrastructures and carries operational metadata that influence client behaviour. Prior work has explored HTTP as a surface for communicating sustainability information. In 2023, Martin proposed an Internet-Draft introducing a \texttt{Carbon-Emissions-Scope-2} response header to communicate emissions associated with request processing \cite{martin-http-carbon-emissions-scope-2-00}. Subsequent discussions in the HTTP working group addressed questions of negotiation, overhead, and field design \cite{httpwg-admin-issue-52-carbon-emissions,ietf-http-wg-mailinglist-carbon-header-thread}. More recently, a 2026 Internet-Draft proposes a structured Sustainability header carrying multiple emission dimensions \cite{besleaga-green-sustainability-header-00}.

However, these efforts primarily address the question of \emph{where} and \emph{what} to disclose at the protocol layer, rather than how the disclosed values can be generated in the first place. In practice, per-request energy and emissions are not directly observable in typical production environments. Hardware-level power telemetry is rarely available at sufficient granularity, and even when it is, attribution is complicated by multi-tenancy, virtualization, varying system load, and request fan-out across internal and external services. As a result, the central challenge is not the design of a disclosure interface, but the construction of a method that can produce meaningful, reproducible per-request estimates under realistic constraints. Without such a method, any protocol-level disclosure remains underspecified, regardless of how the information is encoded or transmitted.

This paper addresses this implementation gap by introducing a practical and
reproducible method for deriving per-request energy and emissions estimates
under realistic deployment constraints. Rather than focusing on the concept of
an HTTP sustainability header, the work concentrates on how such values
can be generated in the first place, given the limited observability of
production systems. The approach deliberately avoids relying on real-time power measurements in
production environments, which are often unavailable, too coarse-grained, or
difficult to interpret. Instead, endpoints are benchmarked offline under
controlled conditions to derive endpoint-specific energy models.
These models make explicit assumptions about system boundaries, measurement
conditions, and allocation choices. At runtime, they can be evaluated using
variables that are already observable at the HTTP boundary, such as request
duration, transferred bytes, or application-level metadata.

HTTP-level disclosure is used in this paper as one concrete instantiation of
this idea, but is not required by the method itself. The core contribution is
the measurement-to-model pipeline that enables per-request sustainability
estimates to be computed, inspected, and reproduced under practical
constraints.

This paper makes four contributions:
(1) It separates the protocol-level question of \emph{what to disclose} from
the systems question of \emph{how to generate request-level values}, and
identifies the latter as a key gap in current HTTP sustainability discussions. (2) It proposes a measurement-to-model pipeline that derives compact,
interpretable, and executable energy models for HTTP endpoints from controlled
offline benchmarks. (3) It demonstrates how these models can be evaluated at the HTTP server boundary and exposed as request-level sustainability metadata, while keeping
underlying assumptions explicit. (4) It evaluates the feasibility, expressiveness, and runtime overhead of the
approach across heterogeneous endpoints, including conventional web services
and AI inference workloads.



%-----------------------------------------------------------------------------------------
\section{Previous Work}
%-----------------------------------------------------------------------------------------
%-----------------------------------------------------------------------------------------

Environmental impacts of digital systems are not merely a property of hardware or data-centre operations, but are also shaped by software design and engineering decisions. Kern et al.\ argue that software and its engineering influence the carbon footprint of information and communication technology (ICT) and therefore belong within the scope of environmental assessment rather than being treated as operational details only \cite{kern2015impacts}. Lago et al.\ further framed sustainability as a property of software quality, thereby moving it into the vocabulary of software architecture and requirements engineering \cite{lago2015framing}. Verdecchia et al.\ later synthesised this perspective under the broader paradigm of Green IT and Green Software, emphasising that software practitioners can influence resource consumption through design, implementation, and operational choices \cite{verdecchia2021green}.

A second strand of research addresses how software-related resource use can be measured and interpreted in a structured way. Guldner et al.\ developed the Green Software Measurement Model (GSMM) as a reference model for assessing the resource and energy efficiency of software products and components \cite{guldner2024gsmm}. Their contribution is particularly relevant here because it makes explicit that software sustainability assessment requires a methodological bridge between raw observations, modelling choices, system boundaries, and reported indicators.

A third body of work demonstrates that digital services delivered over the web can be analysed in environmental terms, even when the underlying infrastructure is complex. Pärssinen et al.\ quantified the environmental impact of online advertising and thereby showed that seemingly lightweight web interactions can aggregate into substantial energy demand and emissions at scale \cite{parssinen2018onlineadvertising}. Sala et al.\ presented Green Web Meter as a real-time digital sustainability monitoring system for websites and web applications, showing that sustainability metrics can be integrated into operational web tooling \cite{sala2024greenwebmeter}. However, these works primarily support assessment and monitoring at the level of websites, services, or traffic aggregates. They do not provide a general mechanism for generating request-specific sustainability values directly at the protocol boundary.

Related work has also shown that web services can adapt their behaviour in response to carbon signals. CASPER, for example, demonstrates that distributed web services can be scheduled and provisioned in a carbon-aware way while respecting service-level constraints \cite{souza2023casper}. At the same time, Fair-CO2 highlights that carbon attribution in cloud environments is intrinsically difficult because workloads share infrastructure and both operational and embodied emissions must be allocated across tenants and services \cite{han2025fairco2}.

Strubell et al.\ helped establish the environmental relevance of machine-learning workloads by quantifying the financial and carbon costs of deep learning in NLP \cite{strubell-etal-2019-energy}. Desislavov et al.\ then shifted attention toward inference and showed that inference energy does not follow simple parameter-count narratives, making deployment-time behaviour an important object of study in its own right \cite{desislavov2023trends}. More recently, Poddar et al.\ benchmarked inference energy across NLP tasks and found strong relationships between energy use, output length, and runtime \cite{poddar2025towards}. 

The Software Carbon Intensity (SCI) specification\cite{gsf-sci}, formalized in ISO/IEC 21031:2024, is an important reference point because it frames software impact as carbon per functional unit \(R\), rather than only as an annual or service-level total. The SCI combines operational emissions from energy use with embodied emissions from the hardware required to run the system.
\[
SCI = \frac{(E \times I) + M}{R}
\]
where \(E\) denotes the energy consumed by the software system, \(I\) the location-based marginal carbon intensity of electricity called grid intensity, \(M\) the embodied emissions allocated to the hardware used by the system, and \(R\) the functional unit.

Existing work has already proposed HTTP-level sustainability disclosure \cite{martin-http-carbon-emissions-scope-2-00,besleaga-green-sustainability-header-00}. These proposals are important because they define possible protocol surfaces for communicating environmental metadata. However, they primarily address what should be communicated at the protocol layer, not how a server can derive credible per-request values under realistic deployment constraints. The practical generation of request-level energy and emissions values therefore remains insufficiently specified.

Taken together, prior work establishes four points. First, sustainability is a legitimate software-quality and software-engineering concern \cite{kern2015impacts,lago2015framing,verdecchia2021green}. Second, credible reporting depends on explicit measurement and modelling choices rather than on raw instrumentation alone \cite{guldner2024gsmm}. Third, web services and AI inference both exhibit environmentally meaningful variation at fine operational granularity. \cite{souza2023casper,sala2024greenwebmeter,poddar2025towards,desislavov2023trends}. Finally adding carbon values to HTTP headers is an ongoing discussion.



%%%% Per-Request Energy Estimation
We model the energy consumption of HTTP endpoints as a function of
request-level features observable at runtime. The goal is to derive compact
models that can be evaluated efficiently at the HTTP boundary while capturing
the dominant factors influencing per-request energy use.

Let \(e\) denote an endpoint and \(x\) denote runtime-observable features. In our setting, these may include request processing time \(t\), request bytes \(b^{req}\), response bytes \(b^{res}\), HTTP method \(m\), response status \(r\), URL-parameter-derived features \(q\), and, where applicable, additional request-scoped metadata \(a\) supplied by the application. We fit a model \(\widehat{E}_e(x)\) that predicts energy per request (in \(\mathrm{mWh}\)).

These features are attractive because they are either already visible to commodity web servers and reverse proxies, or can be exposed with minimal application changes. Request time and transferred bytes capture coarse compute and transfer effects; HTTP method helps distinguish qualitatively different execution paths such as reads, writes, and deletions; response status helps distinguish successful, rejected, redirected, cached, or failing requests; and URL parameters can capture request-specific workload variation such as pagination depth, filter complexity, or user-selected result size.

We use three model families, chosen for interpretability and robust runtime evaluation:

\textbf{Constant model.}
For endpoints whose energy is stable across input variation:
\[
\widehat{E}_e(x) = c_e.
\]

\textbf{Linear model.}
For endpoints expected to scale with execution time, transferred bytes, or categorical request attributes, we use a linear model with suitable encodings:
\[
\widehat{E}_e(x) =
\beta_0
+ \beta_t\, t
+ \beta_{req}\, b^{req}
+ \beta_{res}\, b^{res}
+ \beta_m\, \phi(m)
+ \beta_r\, \psi(r)
+ \beta_q\, \chi(q),
\]
where \(t\) is request processing time, \(b^{req}\) denotes request bytes, \(b^{res}\) denotes response bytes, and \(\phi(m)\), \(\psi(r)\), and \(\chi(q)\) denote encodings of HTTP method, response status, and URL-parameter-derived features, respectively. In practice, \(\beta_0\) can capture constant overheads (e.g., authentication checks), while the remaining coefficients capture scaling effects and systematic differences between request classes. If a deployment does not require method-, status-, or query-specific modelling, these terms can simply be omitted.

\textbf{Curve model (piecewise).}
For endpoints exhibiting non-linear or threshold behaviour (e.g., database writes, batching, or response-size jumps):
\[
\widehat{E}_e(z) = \mathrm{Interp}\bigl(\{(z_i, E_i)\}_{i=1}^k\bigr),
\]
where \(z\) is a selected one-dimensional observable such as request bytes, response bytes, or a numeric URL parameter, and \(\mathrm{Interp}\) is piecewise linear interpolation with defined extrapolation rules (e.g., linear tail).

\textbf{Application-supplied model inputs.}
Some relevant predictors are not normally observable at the HTTP server layer, but are known to the application or to a downstream runtime. In such cases, the application can expose additional request-scoped variables that are then used as inputs to the same modelling approach described above. These variables may include, for example, token count for \texttt{LLM}-style inference, result-set cardinality, number of database records touched, cache-hit indicators from the application layer, prompt length, or other domain-specific workload descriptors.

Formally, such variables can be treated as an extension of the runtime-observable feature vector \(x\). If \(a\) denotes one or more application-supplied variables, the prediction model remains
\[
\widehat{E}_e(x),
\]
with \(x\) now including both server-observable features and application-supplied inputs. For example, for an inference endpoint, token count \(\tau\) may be included as one component of \(x\), yielding a specialised form such as
\[
\widehat{E}_{ai}(\tau) = \alpha_0 + \alpha_\tau \, \tau.
\]
This reflects the empirical observation that many decoder-only models exhibit approximately linear scaling in generated tokens under fixed configuration (temperature, context length, batch size), subject to hardware effects. More generally, however, token count is only one example. The same mechanism can be used for any application-level variable that improves predictive accuracy and can be exposed to the reverse proxy as request-scoped metadata, for example through an internal header or server variable.


\subsection{Operational emissions, embodied emissions, and grid intensity}
\label{subsec:emissions}
We report emissions per request in \(\mathrm{CO_2e}\) by separating operational and embodied components.

\textbf{Operational emissions.}
Given predicted energy \(\widehat{E}_e\) (in \(\mathrm{kWh}\)) and grid intensity \(I\) (in \(\mathrm{gCO_2e/kWh}\)),
\[
\widehat{C}^{op}_e = \widehat{E}_e \cdot I.
\]

\textbf{Embodied emissions as a time-sliced rate.}
Let \(M\) be embodied emissions for the machine (and, if desired, allocated infrastructure) in \(\mathrm{gCO_2e}\). We allocate embodied emissions as a constant rate over an assumed lifetime \(T\) (seconds):
\[
r^{emb} = \frac{M}{T}\ \mathrm{\Bigl(\frac{gCO_2e}{s}\Bigr)}.
\]
For a request with processing time \(t\), embodied emissions are:
\[
\widehat{C}^{emb}_e = r^{emb} \cdot t.
\]
This is an intentionally simple allocation that makes assumptions explicit. Alternative allocations (e.g., utilisation-weighted, capacity-based, or economic allocation) are possible and may be preferable in some settings \cite{iso14040,iso14044,ghgproduct,gsf-sci}.

\textbf{Total per-request emissions.}
\[
\widehat{C}^{req}_e = \widehat{C}^{op}_e + \widehat{C}^{emb}_e.
\]

\subsection{Limitations}
\label{subsec:limitations}
A natural objection is that per-request energy and emissions depend on deployment-specific factors such as hardware generation, power-management settings, co-located workloads, virtualization layers, and regional infrastructure. The same logical endpoint may therefore have different physical footprints in different environments. This objection is valid, but it does not invalidate the usefulness of per-endpoint disclosure. Rather, it clarifies what kind of claim such disclosure can and cannot make.

Our approach does not claim to expose a universally exact physical truth for every request on every deployment. Instead, it provides a reproducible \emph{interface-level estimate} derived from controlled measurements, explicit assumptions, and runtime-observable request features. The purpose of this estimate is not perfect metrological fidelity, but actionable disclosure at the point where software decisions are made. In that respect, per-request sustainability information is similar to other interface-level operational metrics. Latency, throughput, cache-hit ratios, and rate limits are also abstractions over complex underlying systems, yet they remain useful because they are exposed at the boundary where developers and operators can act on them. This abstraction remains useful for four reasons.

First, \emph{software interfaces require decision-oriented metrics rather than hardware-complete models}. A client deciding whether to retry a request, reduce a prompt, change a pagination limit, or call a different endpoint does not primarily need a full physical simulation of the server environment. It needs a signal that is stable enough to differentiate request classes and reflect the relative effect of its choices. For such decisions, comparative sensitivity is often more important than absolute physical precision.

Second, \emph{the model is explicit and therefore inspectable}. The disclosed value is not presented as an unqualified sensor reading. It is the output of a documented model with a defined system boundary, a stated measurement setup, and named inputs. This makes assumptions contestable and reproducible. In sustainability assessment, such transparency is often more valuable than apparent precision from opaque instrumentation.

Third, \emph{deployment dependence can be managed through calibration rather than treated as a fatal flaw}. The model structure can remain stable even when coefficients differ across machines or hosting environments. Providers can therefore calibrate the same model family to their own infrastructure, publish the reference environment, or provide deployment-specific registries. In this sense, hardware variance affects parameter values, but not the viability of the disclosure mechanism itself.

Fourth, \emph{silence is often worse than approximate but qualified disclosure}. If the only acceptable value were a perfectly attributable real-time physical measurement, then most contemporary HTTP services could disclose nothing at all. Multi-tenancy, cloud opacity, and downstream fan-out would make request-level reporting practically impossible. A qualified estimate with explicit assumptions and uncertainty is therefore preferable to the absence of any per-request information.

The key point is that request-level disclosure should be understood as an \emph{operational sustainability interface}, not as a claim of hardware-independent exactness. Its value lies in making environmental consequences visible in a form that software can use, compare, and optimise against. Hardware variance does not remove that value; it simply means that the disclosure must be interpreted as model-based, scoped, and context-dependent.

